\label{SubSec:HeatVulnerabilityModel}

\paragraph{Impact of temperature on labour productivity\\}
\textbf{\\}
The heat vulnerability model presented in this section is based on the approach introduced in \cite{NeidellEtAl:2021}. The paper uses survey data to estimate labour allocation decisions in the United States (US) based on temperature. It does not extend the analysis beyond the US. It replicates previous research done (original GZN method \cite{NeidellEtAl:2014}.) whilst extending the period of data used and adding an assessment based on the economic cycle: the main innovation is that it includes the economic cycle by splitting out the 2008 financial crisis first through segmented regressions and following them by using an indicator variable: the non-recession period is 2003-2007 and 2015-2018. The `Great Recession' period is 2008-2014. The methodology in only applied to climate exposed sectors: agriculture, forestry, fishing and hunting, construction, mining, and transportation and utilities.

The paper's main conclusions are:
\begin{itemize}
    \item A statistically significant impact of 2.6 minutes lost per degree of temperature above $90$°F during normal economic periods and no relationship during a recession. This result is converted into the Celsius scale as \emph{physrisk} decided to take that scale as a reference. The 2.6 minutes is multiplied by a scaling factor of 1.8, which returns \textbf{an impact of 4.7 minutes lost per degree of temperature above $\mathbf{32.2}$°C}.
    \item When using an indicator variable and linear regression the estimated impact was 5.6 minutes under Fahrenheit scale (respectively 10.08 minutes under Celsius scale) during normal economic periods, but the parameter was not significant at the 90\% level.
    \item No relationship between temperature and work allocation with temperatures below  $32.2$°C.
    \item Focus on labour allocation decisions, it does not account for other impacts such as reducing productivity.
\end{itemize}
While the results are relevant and significant, is it important to highlight the following disclaimers on the results:
\begin{itemize}
    \item There are some questions of the reliability of the forecasts in the long run as climate change will likely result in structural changes: on a long-term basis, adaptive solutions might be considered by people to adjust their work productivity with respect to temperature rise.
    \item The conclusion holds for US (so maybe also for the EU as well other developed countries), but not for developing countries which experience more economic turmoil periods.
\end{itemize}
The paper attempts to estimate economic cost (assuming the impact of 4.7 minutes lost per degree of temperature above $32.2$°C during normal economic periods) however it only focuses on the direct costs and does not account for feedback effects (reducing the labour productivity will result in a decrease of the products available, wages, demand, etc.). Figure \ref{fig:economiccost} provides the results, as extracted from the paper:
\begin{figure}[h]
    \centering
    \includegraphics[scale = 0.6]{plots/economic cost.png}
    \caption{Economic cost}
    \label{fig:economiccost}
\end{figure}
Last but not least, it is worth noting that the results of this paper are interpreted as the impact of \textbf{chronic increase in temperature}.
%NEW
%%%%%%%%%%%%%%%%%%%%%%%%%%%%%%%%%%%%%%%%%%%%%
The GZN hazard  model is based on the projections of the climate variable `daily maximum temperature'. The projected data spans for 20 years. Then, the following statistical processing is used to compute the daily cooling degree days: if the daily number of degrees is higher than $32.2$°C then the cooling degree days is equal to the number of degrees, and 0 if not. Once the 20-year projected time series of daily cooling degree days is computed, the indicator of the GZN hazard data is calculated as an annual average of the cooling degree days (over all the days within the 20-year period). That indicator will be used afterwards as an input in the impact function to compute the number of minutes of labour productivity loss. %Figure \ref{fig:GZL-Hazard} provides a summary of the methodology:

$T^\text{max}$ is the daily maximum near-surface temperature. This is a point of attention, given that heating and cooling degree day indicators are often calculated using the daily \emph{average} near-surface temperature.
\begin{equation}
    \label{Eq:degree_days}
    I^\text{dd} = \frac{365}{n_y} \sum_{i = 1}^{n_y} |  T^\text{max}_i - T^\text{ref} |
\end{equation}

\paragraph{Uncertainty around the vulnerability Heat model\\}
\textbf{\\Overview}

The assessment is based on the research performed by Zhang and Shindel, which reviews the uncertainty in the heat risk literature \cite{ZhangAndShindell:2021}. This paper provides context around the uncertainty that exists in the result discussed in previous section, which is mainly explained by the methodology used for the estimation of the impact of temperature on labour productivity.
\begin{itemize}
    \item \cite{ZhangAndShindell:2021} provides an analysis of the \textbf{differential forecasts} between using the \textbf{GZN method} (as reference to Graff-Zivin and Neidell) documented in \cite{NeidellEtAl:2021}, versus the WetBulb Globe Temperature methodology (\textbf{WBGT method}) which includes other climate factors in addition to temperature: humidity, wind speed and heat radiation -- figure \ref{fig:WBGT} \footnote{Extracted from \cite{ZhangAndShindell:2021} -- p.4} provides the WBGT detailed approach.
    \item Another major source of differentiation is that the GZN method focuses on changes in labour allocation decisions while the WBGT method focuses on the physiological impacts of rising temperatures.
\end{itemize}
\begin{figure}[H]
    \centering
    \includegraphics{plots/WBGT.png}
    \caption{WBGT method}
    \label{fig:WBGT}
\end{figure}
The $\alpha_1$ and $\alpha_2$ are the parameters derived for the three different work intensities: light (24.64, 22.72), medium (32.98, 17.81), and heavy (30.94, 16.64) work.
\begin{itemize}
    \item Light work includes all service sectors, such as trade, all retail sales, wholesale trade and commission trade, hotels and restaurants, repairs of motor vehicles and personal and household goods, retail sale of automotive fuel, post and telecom, financial services, insurance, recreational and service activities, public administration, dwellings real estate, and other businesses.
    \item Medium work includes food, textile and wood, machinery and electronic equipment, and other industries and transport sectors.
    \item Heavy work includes agriculture, forestry and fishery, extraction sectors, and construction sectors. These are the sectors considered in the GZN model methodology, which are the climate exposed sectors. Hence in this study, $\alpha_1$ and $\alpha_2$ of the heavy work intensity are applied $(30.94, 16.64)$.
\end{itemize}

Note that there are differences in the functional forms applied in the original GZN method \cite{NeidellEtAl:2014} and the approach presented in \cite{NeidellEtAl:2021}, with original GZN method using one linear regression with dummy variables for temperature buckets, while \cite{NeidellEtAl:2021} uses multiple linear regressions with one variable reflecting the breach of the maximum temperature around anchor points (less than $70$°F, $90$°F, $90$°F and above). The WBGT methodology uses a non-linear function to relate labour loss to the WBGT consolidated measure.

Figure \ref{fig:CostsByRCP} \footnote{Extracted from \cite{ZhangAndShindell:2021} -- p.11} provides the forecasts of labour lost millions of 2016 USD (constant USD value). Most notably the original GZN produces more optimistic forecasts of the cost of labour lost (lower) than the WBGT method. Note that RCP8.5\footnote{The Intergovernmental Panel on Climate Chance modelling are based on representative concentration pathways (RCPs), which represent different emissions projections under basic, plausible economic and social assumptions, while staying within physical constraints. RPCs are constructed by back-calculating the amount of emissions that would result in a given amount of radiative forcing (which is the difference between solar radiation (energy) absorbed by the Earth and energy radiated back into the space) that would then result in a given amount of warming} refers to the scenario of high emissions and RCP4.5 refers to the scenario of moderate emissions.
\begin{figure}[h]
    \centering
    \includegraphics{plots/CostsByRCP.png}
    \caption{Costs By RCP under GLZ method versus WBGT method}
    \label{fig:CostsByRCP}
\end{figure}

\textbf{Detailed approach: vulnerability around GLZ model}

There are two identified areas where uncertainty exists in the forecasts in \cite{NeidellEtAl:2021}: the economic cycle and the model parameter uncertainty.

The \textbf{economic cycle} is one area where uncertainty exists in the forecasts.
The paper shows that labour allocation decisions are sensitive to where in the economic cycle the US is; during a recession there does not appear to be a relationship between labour allocation and temperature. In order to measure the uncertainty explained by the economic cycle, one might consider the probability of a recession as a Bernoulli random variable with a probability p. Based on this there are two possibly approaches. A first approach is a Monte-Carlo like approach where one can randomly sample 1 or 0 depending on whether a recession occurs or not at each period on a time path. A second approach would be to use the expected value of the probability of the recession and estimate the impact as:
\begin{equation}
    \label{Eq:economiccycle}
        \text{forecast minutes lost} = p \times 0 + (1-p) \times \text{EV}\left(X\right)
\end{equation}
Where $X$ is the variable that refers to Minutes Lost during normal economic cycle
The second approach is attractive in its simplicity and ensuring the model does not lose its focus.
One concern is that the probability of recession, in reality, may be related to the realisation of climate related risks. Hence, there might be a over/under estimation of the probability unless the impact of the climate risk is also considered. Historical model estimated recession probabilities can be sourced from \cite{ChauvetEtAl:2024}. However, we did not go further in the measurement of the uncertainty explained by the economic cycle because it requires us to focus on the modelling of another parameter ($p$, the probability of the recession to occur), which is not the purpose of this work.

%%%%%%%%%%%%%%%%%%%%%%%%%%%%%%%%%%%%%%%%%%%%
% to add in the references: https://github.com/os-climate/physrisk/blob/main/methodology/PhysicalRiskMethodologyBibliography.bib
%%%%%%%%%%%%%%%%%%%%%%%%%%%%%%%%%%%%%%%%%%%%
%@article{SmoothedU.S.RecessionProbabilities:2022,
%  title={Smoothed U.S. Recession Probabilities},
%  author={Jeremy Piger},
%  publisher={University of Oregon}
%  year={2022},
%}
Instead, we focus on the \textbf{model parameter uncertainty approach}. In \cite{ZhangAndShindell:2021}, there is a linear relationship between temperature and work minutes lost ($\beta$),and a constraint is applied to ensure that total time allocation sums to 24 hours, which returns a non-linear regression model at the end. Given that the main coefficient of interest is denoted $\beta$, an inference is applied assuming that $\beta$ follows a Student's-t distribution:
\begin{equation}
 \label{Eq:uncertaintyStudentT}
 \delta \beta \sim T_{\left(\beta, \mathrm{SE}, N-K\right)}
\end{equation}
Where $\text{SE}$ is the Standard Error of the coefficient and $(N-K)$ is the number of degrees of freedom, $N$ is the number of observations and $K$ is the number of model parameters. Hence, one can estimate the coefficients of the confidence interval (CI) at a given level of confidence CI\%:
\begin{equation}
 \label{Eq:CIStudent}
 \beta_{\mathrm{CI}\%} = \beta \pm T(p) \times \text{SE}
\end{equation}
Where $T(p)$ donates the probability density function (pdf) of the student $T$ distribution with probability $p$ which corresponds to the CI\%, the standard error $\text{SE}=2.23$ minutes and $\beta=-4.68$.
%%%%%%%%%%%%%%%%%%%%%%%%%%%%%%%%%%%
%START HERE
%%%%%%%%%%%%%%%%%%%%%%%%%%%%%%%%%%%
Figure \ref{fig:GZL-Vulnerability} shows the uncertainly around the GZN vulnerability model, based on the assumptions above, at $95\%$ confidence level\footnote{In order to compute the figure, N is taken from the paper \cite{TemperatureAndWork:2021}: N 11,732, and K is assumed to be 0 as it is very small relatively to N. One could do further researches to check the exact number of K, but that will not impact the results. K includes the number of control variables which provide information on demographic data -- age, gender, education, employment status, income, etc. -- and climate variables such as level of precipitation and snow.}.
\begin{figure}[h]
    \centering
    \includegraphics[scale = 0.8]{plots/GZL-Vulnerability.png}
    \caption{Estimate of daily minutes of labour loss}
    \label{fig:GZL-Vulnerability}
\end{figure}
As the degrees of freedom increases, the t distributions converge to the standard normal. Therefore, for simplicity reasons in the context of this work, it is assumed that that the daily labour productivity impact $\beta$ can be measured as:
\begin{equation}
    \label{Eq:uncertainty1}
        \beta \sim \mathcal{N}_{\left(m,\text{SE}^{2}\right)}
\end{equation}
where $m = 4.68$ minutes and $\text{SE}=2.23$ minutes. As an example, consider an $1.5$°C degree day increase in the temperature. We can then multiply through the normal distribution as shown below:
%The buckets of the daily maximum temperature above $32.2$°C  are defined as an incremental increase of $1.5^\circ %C$, starting from $0$°C to $18$°C: $\begin{pmatrix} 0 & 1.5 & 3 & ... & 18 \end{pmatrix}$. For a $1.5$°C daily temperature increase, the uncertainty around the daily labour productivity impact is given by the following %normal distribution:
\begin{equation}
    \label{Eq:uncertainty2}
    1.5 \times \beta \sim \mathcal{N}_{\left(1.5 \times \beta,(1.5 \times \text{SE})^{2}\right)}
\end{equation}
This can be generalised for a degree day increase of x using the following distribution:
\begin{equation}
    \label{Eq:uncertainty3}
    x \times \beta \sim \mathcal{N}_{\left(x \times \beta,\left(x\times\text{SE}\right)^{2}\right)}
\end{equation}
For example, there is 1\% risk that the lost labour productivity exceeds $29.6$ minutes per day if the maximum daily temperature exceeds at $32.2$°C by $3$°C. This number is computed as
$$
\mathcal{N}^{-1}_{\left(m',\left(\text{SE}'\right)^2\right)}\left(1\%\right)
$$
where $m'= 3 \times m = -14.013$ minutes and $\text{SE}'= 3 \times \text{SE} = 6.69$ minutes. The $99\%$ confidence interval of the lost labour productivity per day if the maximum daily temperature exceeds at $32.2$°C by $3$°C is $\left[+1.6\text{ min}, -29.6\text{ min}\right]$.
%Adding point regarding impact as a percentage of total labour
The estimated loss of labour time is then transformed into a percentage estimate:
\begin{equation}
    \label{Eq:uncertainty3}
    \text{estimated loss of labour time (\%)}=  \frac{\text{labour lost}}{\text{total minutes worked in a year}}
\end{equation}
For this the figures reported by the Organisation for Economic Co-operation and Development (OECD) are used, specifically the 2021 estimate for the United States of America. This is to ensure the alignment with the region where the labour parameters are estimated in. Buckets of labour lost are defined within the range of no impact to 100\%, with estimated marginal probabilities for each bucket returned.

The main weakness of the Temperature and Work article (which is more likely to be close to the original GZN method than to the WBGT method) is that it does not take into account of higher levels of optimism in the original GZN method which was noted in figure \ref{fig:CostsByRCP}. Hence, the results might be underestimating the actual impact. The next paragraph explores the vulnerability model around the WBGT method and its uncertainty.

\textbf{Detailed approach: vulnerability around WBGT model}

%NEW
%%%%%%%%%%%%%%%%%%%%%%%%%%%%%%%%%%%%%%%%%%%%%
The WBGT hazard model is based on the projections of the climate variables 'daily temperature' and 'daily humidity'. The projected data spans for 20 years. Then, the computations provided in figure \ref{fig:WBGT} are used to compute the daily WBGT, which is considered as the indicator of the WBGT hazard data that will be afterwards considered as an input in the impact function to get the daily Work Ability (WA) projections over the 20-year period. Finally, the impact deriving from the WBGT model is computed as the annual average of the daily projected WA.
%\begin{figure}[h]
%    \centering
%    \includegraphics[scale = 0.8]{plots/WBGT-Hazard.PNG}
%    \caption{WBGT and labour availability}
%    \label{fig:WBGT-Hazard}
%\end{figure}
Note that another indicator deriving from the WBGT hazard model is the Work Loss (WL), which is computed as $\text{WL} = 1 - \text{WA}$. It is another way to do the assessment leading to the same results.

The aggregation of the outputs of the GZN vulnerability model (minutes of productivity labour loss) and the WBGT vulnerability model (WA), returns the effective number of working hours, which represents a way to measure the uncertainty around the GZN vulnerability model. %Figure \ref{fig:Aggregation} provides the aggregation process and results:
%\begin{figure}[h]
%    \centering
%    \includegraphics[scale = 0.8]{plots/Aggregated Hazard WBGT GZN.PNG}
%    \caption{Aggregation of GZN model and WBGT model}
%    \label{fig:Aggregation}
%\end{figure}
If the WL was used instead of the WA, then it will be multiplied by the hours worked derived from the GZN model to get the annual total labor loss due to heat.

Given the modelling assumptions around the parameters used to compute the WA in the WBGT model, it is important to measure the uncertainty around these parameters, $\alpha_1$ and $\alpha_2$, which depend on the work intensities.

The source paper provides the parameters $\alpha_1$ and $\alpha_2$ for three different work intensities low, medium and high with industries mapped to each intensity. These categories are broad and do not account for variance within and industry and between industries in the same category. To account for this uncertainty the WBGT approach was adjusted to include uncertainty around the industry.

Consider an asset which is market as in a high risk sector. We assume that the work ability is uniformly distributed with a mean equal to the $\text{WA}_H$ and a floor (a) and ceiling (b) equidistant from the mean. We assume that the floor a is halfway between $\text{WA}_H$ and $\text{WA}_M$. So a and b can be estimated based on the below formulae:
\begin{equation}
    \label{Eq:WBGT_Floor}
     a = \text{WA}_H - \frac{\text{WA}_H - \text{WA}_M}{2}\,
\end{equation}
\begin{equation}
    \label{Eq:WBGT_Ceil}
     b = \text{WA}_H + \frac{\text{WA}_H - \text{WA}_M}{2}\,
\end{equation}
And the WBGT work ability can be represented based on the below formula:
\begin{equation}
    \label{Eq:WBGT_Uniform}
     \text{WA} \sim \mathcal{U}(a ,b)\,
\end{equation}
With this we can estimate the variance of the WBGT work ability using the standard formula for the variance of a Uniform distribution:
\begin{equation}
    \label{Eq:WBGT_Variance}
     Var\left(\text{WBGT}\right) = \frac{(a-b)^2}{ 12}
\end{equation}
In order to get to a final work ability we multiple the output of the GZN model by the ouput of the WBGT model.
\begin{equation}
    \label{Eq:Final_Mean_WA}
     \text{effective work} = \left(1- \text{estimated loss of labour time}\right) \times \text{WA}_H
\end{equation}
As we have uncertainty both for the GZN component as well as the WBGT component we need to estimate a joint variance. We assume that the epistemic uncertainty in the GZN model is uncorrelated with the variance in the sector uncertainty of the work ability measure. To estimate the variance of the product of both the GZN and WBGT components we using the following formula:
\begin{equation}
\begin{aligned}
    \label{Eq:Var_Joint}
     Var\left(\text{effective work}\right) = \left(1- \text{estimated loss of labour time}\right)^2 \times Var\left(\text{WBGT}\right) + \\
      \text{WA}_H^2 \times Var\left(\text{GZN}\right) +  Var\left(\text{WBGT}\right) \times Var\left(\text{GZN}\right)
      \end{aligned}
\end{equation}
To get the final impact we assume that the product of the two variables are normally distributed. The final work ability distribution can be represented as below:
\begin{equation}
\begin{aligned}
    \label{Eq:Final_Result_Heat}
     \mathcal{N}_{\left(\text{effective work}, Var\left(\text{effective work}\right)\right)}
      \end{aligned}
\end{equation}
\setcounter{secnumdepth}{3}